% =============================================================================
% Monthly 4D-Var estimation of terrestrial CO2 fluxes in GCHP: science report
% Base draft for a paper / science report.  Companion documents:
%   monthly_sigma_gradient.tex  - derivation of the monthly adjoint gradient
%   monthly_4dvar_software.tex   - software architecture and workflow
% =============================================================================
\documentclass[12pt]{article}

% Latin Modern ships Type 1 outlines; plain Computer Modern needs MetaFont (mf),
% which is absent on the NAS hosts and aborts the build.  Do not add T1 fontenc
% or the bm package for the same reason.
\usepackage{lmodern}
\usepackage{amsmath,amssymb,amsfonts}
\usepackage{geometry}
\usepackage{graphicx}
\usepackage{booktabs}
\usepackage{hyperref}
\usepackage{xcolor}
\newcommand{\bm}[1]{\mathbf{#1}}

\geometry{margin=1in}

\newcommand{\xco}{X\mathrm{CO}_2}
\newcommand{\co}{\mathrm{CO}_2}
\newcommand{\Fprior}{F^{\mathrm{p}}}
\newcommand{\tstart}{t_{\mathrm{start}}}
\newcommand{\tend}{t_{\mathrm{end}}}

\title{\textbf{Monthly 4D-Var Estimation of Terrestrial CO$_2$ Fluxes
with GEOS-Chem High Performance}\\
\large An observing-system simulation experiment}
\author{}
\date{\today}

\begin{document}
\maketitle

\begin{abstract}
We describe a four-dimensional variational (4D-Var) system for estimating
surface CO$_2$ fluxes from column-averaged dry-air mole fraction ($\xco$)
observations, built on the adjoint of GEOS-Chem High Performance (GCHP). The
control vector is a set of monthly multiplicative scaling fields applied to the
prior terrestrial ecosystem respiration (TER) flux, so the analysis resolves the
seasonal evolution of the biospheric flux adjustment rather than a single
annual correction. The system is exercised in an observing-system simulation
experiment (OSSE): synthetic $\xco$ retrievals are sampled from a known ``true''
flux field along satellite ground tracks, and the assimilation attempts to
recover the truth from a biased prior. This report sets out the scientific
formulation---cost function, control variables, observation operator, and prior
error model---and the adjoint-based gradient that makes the monthly extension
tractable. We summarise diagnostics from a single-month proof-of-concept
assimilation, identify the two principal obstacles that the diagnostics expose
(a noise-dominated adjoint gradient and an under-weighted background term), and
lay out the monthly, full-year experiment and the analysis plan. The companion
note \texttt{monthly\_sigma\_gradient.tex} derives the monthly gradient; the
companion note \texttt{monthly\_4dvar\_software.tex} documents the
implementation.
\end{abstract}

\tableofcontents

% =============================================================================
\section{Introduction and motivation}
% =============================================================================

Atmospheric measurements of $\co$ constrain surface fluxes through inverse
(``top-down'') estimation: given a chemical transport model, a prior flux
estimate, and observations of atmospheric $\co$, one seeks the flux field most
consistent with both the prior and the observations. The terrestrial biosphere
is the dominant source of interannual and seasonal variability in the
atmospheric $\co$ growth rate, and its net exchange---the small residual of two
large opposing fluxes, photosynthetic uptake and respiration---is the least
well-constrained term in the global carbon budget. Space-borne $\xco$
retrievals (e.g.\ OCO-2) sample this signal densely enough to constrain fluxes
at sub-continental scales, but only if the inversion resolves the flux in time:
a single annual scaling cannot represent the phase and amplitude of the
biospheric seasonal cycle, which is precisely the quantity of scientific
interest.

This motivates a control vector with monthly temporal resolution. We estimate,
for each calendar month $m$ and each surface grid cell $x$, a multiplicative
scaling factor $\sigma_m(x)$ on the prior TER flux. Twelve monthly fields over a
one-year window let the analysis move the phase and amplitude of the terrestrial
flux independently in each month while keeping the control vector small enough
for a gradient-based optimiser. The remaining flux components---fossil-fuel
combustion, air--sea exchange, and the gross photosynthetic and net biosphere
terms---are held at their prior values, so the experiment isolates the
respiration adjustment that the observations can support.

We evaluate the system in an OSSE rather than with real retrievals. In an OSSE
the ``true'' fluxes are known by construction, so the analysis error---the
difference between the recovered and true fluxes---is available exactly, which
is impossible with real data. This makes the OSSE the natural setting in which
to validate the assimilation machinery, quantify the information content of the
observing system, and expose pathologies of the optimiser and the error model
before committing to a real-data campaign.

% =============================================================================
\section{The forward model}
% =============================================================================

The forward model is a GCHP simulation of atmospheric $\co$ as a single passive
tracer transported by assimilated meteorology. Surface fluxes enter through the
Harmonized Emissions Component (HEMCO), which assembles the net flux from
independently specified components:
%
\begin{equation}
  F_{\mathrm{total}}(x,s) = F_{\mathrm{ff}} + F_{\mathrm{ocean}}
     + F_{\mathrm{GPP}} + F_{\mathrm{NBE}}
     + \sigma_{m(s)}(x)\,\Fprior_{\mathrm{TER}}(x,s),
  \label{eq:fluxsum}
\end{equation}
%
where $m(s)$ is the calendar month containing time $s$. The five terms are
fossil fuel ($F_{\mathrm{ff}}$), air--sea exchange from an ocean biogeochemistry
product ($F_{\mathrm{ocean}}$), gross primary production ($F_{\mathrm{GPP}}$),
net biosphere exchange from a data-assimilation product ($F_{\mathrm{NBE}}$),
and terrestrial ecosystem respiration ($\Fprior_{\mathrm{TER}}$). Only the last
carries the control variable $\sigma$; the seasonal and spatial structure of the
prior TER flux is retained, and $\sigma$ rescales it.

The simulation is run at C24 cubed-sphere horizontal resolution with a
$30$-minute physics timestep. The control fields $\sigma_m$ are defined on a
$72\times46$ regular latitude--longitude grid and mapped to the cubed sphere by
HEMCO. A single forward integration over the assimilation window produces
simulated $\xco$ along the satellite ground tracks, which is the model
counterpart of the observations.

% =============================================================================
\section{Observations and the OSSE design}
% =============================================================================

\subsection{Observation operator}

The observations are $\xco$ soundings located along satellite ground tracks. The
observation operator $H$ maps the three-dimensional model $\co$ field to a
column-averaged dry-air mole fraction at each sounding location and time by a
pressure-weighted vertical integral. In the forward run, the model is sampled at
the sounding space--time coordinates through a satellite-track diagnostic, so
$H(\bm{x})$ is available for every observation without a separate interpolation
step.

\subsection{Synthetic truth and observations}

The OSSE ``truth'' is a forward simulation driven by a set of terrestrial fluxes
distinct from the prior, so that the true monthly scaling differs from unity in a
known way. Synthetic $\xco$ observations are generated by sampling the truth run
along the real satellite tracks and adding representative retrieval noise. The
assimilation is then initialised from the prior ($\sigma_m \equiv 1$) and asked
to recover the truth. Because the truth is known, the monthly analysis increment
$\sigma_m - 1$ can be compared directly against the true scaling, month by month,
which is the central evaluation of the experiment.

% =============================================================================
\section{Variational formulation}
% =============================================================================

\subsection{Cost function}

Let $\bm{\sigma}$ denote the stacked control vector of all monthly fields,
$\bm{\sigma} = (\sigma_1,\dots,\sigma_{12})$, each $\sigma_m$ a field on the
control grid. The analysis minimises
%
\begin{equation}
  J(\bm{\sigma}) \;=\; J_{\mathrm{obs}}(\bm{\sigma}) + J_{\mathrm{b}}(\bm{\sigma}),
  \label{eq:cost}
\end{equation}
%
with an observation term and a background (prior) term. The observation term is
the $\xco$ misfit summed over all soundings $i$ in the window,
%
\begin{equation}
  J_{\mathrm{obs}}(\bm{\sigma}) \;=\; \tfrac{1}{2}
    \sum_i \frac{\big(H_i(\bm{\sigma}) - y_i\big)^2}{\varepsilon_i^2},
  \label{eq:jobs}
\end{equation}
%
where $y_i$ is the observed $\xco$, $H_i(\bm{\sigma})$ the simulated $\xco$ at
that sounding, and $\varepsilon_i$ the observation-error standard deviation
(measurement plus representation error). The background term penalises departures
of each monthly field from the prior,
%
\begin{equation}
  J_{\mathrm{b}}(\bm{\sigma}) \;=\; \frac{1}{2\sigma_b^2}
    \sum_{m=1}^{12} \lVert \sigma_m - 1 \rVert^2,
  \label{eq:jb}
\end{equation}
%
a diagonal background-error model with a single scalar standard deviation
$\sigma_b$ common to all months and cells. Section~\ref{sec:challenges}
discusses the limitations of this diagonal choice and the planned replacement by
a spatially correlated background.

\subsection{Control variable and its gradient}

The minimisation is performed with a limited-memory quasi-Newton method
(L-BFGS-B), which requires the gradient $\nabla_{\bm{\sigma}} J$. The background
gradient is immediate from \eqref{eq:jb},
%
\begin{equation}
  \nabla_{\sigma_m} J_{\mathrm{b}} = (\sigma_m - 1)/\sigma_b^2 .
\end{equation}
%
The observation gradient is supplied by the GCHP adjoint. The key result, derived
in the companion note \texttt{monthly\_sigma\_gradient.tex} and summarised here,
is that the archived adjoint surface diagnostic $A(x,t) \equiv
\texttt{SurfaceFluxAdj\_CO2}$ is a \emph{backward accumulator in emission time},
%
\begin{equation}
  A(x,t) = \sum_{s\,\ge\,t} \lambda(x,s)\,\Fprior_{\mathrm{TER}}(x,s)\,\Delta t,
  \qquad A(x,\tend)=0,
  \label{eq:accum}
\end{equation}
%
where $\lambda(x,s) = \partial J_{\mathrm{obs}}/\partial F(x,s)$ is the
sensitivity to an instantaneous flux input. Because $\sigma_m$ multiplies the
prior flux only within month $m=[t_m,t_{m+1})$, the monthly observation gradient
is the \emph{difference} of the accumulator across consecutive month boundaries,
%
\begin{equation}
  \boxed{\;g_m(x) \;\equiv\; \frac{\partial J_{\mathrm{obs}}}{\partial \sigma_m(x)}
    \;=\; A(x,t_m) - A(x,t_{m+1}). \;}
  \label{eq:gm}
\end{equation}
%
Two properties of \eqref{eq:gm} are used as runtime consistency checks: the
terminal condition $A(x,\tend)=0$, and the telescoping identity $\sum_m g_m =
A(x,\tstart)$, i.e.\ the twelve monthly gradients must sum to the single
whole-window gradient. Critically, the adjoint accumulator is weighted by the
\emph{prior} TER flux (this is verified in the adjoint source and configuration),
so no $1/\sigma_m$ renormalisation of \eqref{eq:gm} is required, and the adjoint
run never needs the monthly scaling. A corollary is that only the forward model
carries the monthly plumbing.

\subsection{A note on temporal information flow}

It is worth stating explicitly, because it is a common source of confusion, that
the monthly differencing \eqref{eq:gm} does \emph{not} discard the ability of a
late-window observation to constrain an early-window flux. A December $\xco$
observation informs January fluxes through the adjoint transport, which is
encoded inside $\lambda(x,s)$ (a sum over all observations at times $\ge s$).
Differencing restricts only the \emph{emission-time} range to month $m$, because
$\sigma_m$ multiplies no flux outside its own month. The two sums are indexed by
different variables and are not in tension; see the companion note.

% =============================================================================
\section{Experimental configuration}
% =============================================================================

\begin{center}
\begin{tabular}{@{}ll@{}}
\toprule
Setting & Value \\
\midrule
Transport model & GEOS-Chem High Performance (GCHP), $\co$ tracer \\
Horizontal resolution & C24 cubed sphere \\
Physics timestep & $1800$ s (chemistry and dynamics) \\
Assimilation window & 2016-01-01 to 2017-01-01 (12 months) \\
Control vector & $\sigma_m(x)$, $m=1,\dots,12$, on a $72\times46$ lat--lon grid \\
Scaled flux component & terrestrial ecosystem respiration (TER) only \\
Prior / first guess & $\sigma_m \equiv 1$ (unperturbed prior TER) \\
Background error $\sigma_b$ & $0.2$ (uncorrelated, single scalar) \\
Observations & synthetic $\xco$ along satellite ground tracks (OSSE) \\
Optimiser & L-BFGS-B, limited memory, unit step \\
\bottomrule
\end{tabular}
\end{center}

\noindent
Each outer iteration comprises a forward integration over the window, evaluation
of $J_{\mathrm{obs}}$ and the observation forcing, a backward adjoint integration
that produces $A(x,t)$, and one L-BFGS-B update. The monthly gradient
\eqref{eq:gm} is formed by reading $A$ at the thirteen month boundaries.

% =============================================================================
\section{Proof-of-concept diagnostics}
% =============================================================================

Before the full monthly, full-year experiment, the machinery was exercised in a
reduced configuration with a single scaling field over a one-month window, which
shares every software component with the monthly system except the temporal
differencing. Over thirteen outer iterations the observation cost decreased
monotonically (after a first-step overshoot discussed below), and the recovered
scaling field remained physically plausible, staying within roughly
$[0.5,1.6]$. These runs establish that the forward/adjoint/optimiser loop is
functional and that the gradient drives $J$ downward. They also expose two
problems that must be addressed for a credible monthly analysis.

% =============================================================================
\section{Principal challenges exposed by the diagnostics}
\label{sec:challenges}
% =============================================================================

\subsection{A noise-dominated adjoint gradient}

The most consequential finding is that the adjoint gradient field is dominated by
small-scale noise. Between consecutive outer iterations, over which the control
field moves only slightly and $J_{\mathrm{obs}}$ changes by a fraction of a
percent, the gradient field decorrelates almost completely, its root-mean-square
magnitude varies by more than an order of magnitude, and isolated cells
occasionally spike to many times the typical value. The persistent hotspots
coincide with regions of deep tropical convection.

A dedicated determinism test rules out a code defect as the cause: rerunning the
same adjoint integration twice from byte-identical inputs produced
bitwise-identical output. The gradient noise is therefore a genuine feature of
the adjoint solution---chaotic amplification of tiny differences by the adjoint
of moist convective transport---rather than a nondeterministic bug. This is a
known difficulty for adjoints of atmospheric transport with active convective
parameterisations. Its practical effect is that the optimisation has behaved
more like noisy gradient descent stabilised by the background term than like a
clean quasi-Newton minimisation, and it implies that the monthly differencing
\eqref{eq:gm}, being a difference of two noisy fields, will have a poorer
signal-to-noise ratio still. Mitigation is discussed in
Section~\ref{sec:outlook}.

\subsection{An under-weighted background term}

In the proof-of-concept runs the background term contributed a negligible
fraction of the total cost, and its gradient was one to two orders of magnitude
weaker than the observation gradient at the most active cells. The prior
therefore constrained the solution only weakly where the noisy observation
gradient was largest, allowing the analysis to chase noise in those cells. The
root cause is the diagonal, uncorrelated specification of both the
observation-error covariance (which overcounts information from spatially
correlated retrieval errors) and the background-error covariance. As an interim
measure the background standard deviation was tightened; the durable fix is a
spatially correlated background, described next.

\subsection{Finite-window edge effects}

Fluxes emitted late in the window have little time to influence any observation
before the window closes, so their gradients are intrinsically small and the
corresponding months remain close to the prior. This is a property of any
finite-window inversion under a correct gradient, not an artefact of the monthly
formulation; if the late months are of scientific interest, the window can be
extended past the last month to be optimised.

% =============================================================================
\section{The monthly, full-year experiment}
% =============================================================================

The monthly system extends the proof of concept in three ways: the control
vector carries twelve monthly fields instead of one; the forward model reads a
distinct scaling field per month through a time-varying input; and the gradient
is formed by the boundary differencing \eqref{eq:gm}. The monthly plumbing has
been implemented and unit-tested, and its correctness at the level of flux
application (the right monthly field is applied to TER in the right month, and no
other flux component is affected) has been verified from a dedicated forward
test run. A finite-difference test of the gradient convention---confirming that
the difference \eqref{eq:gm}, and not the raw accumulator value, matches a
brute-force perturbation of the cost function---is the gating check before the
full multi-month optimisation, and is run over a short window at reduced cost. A
reduced, three-month end-to-end run serves as an integration test of the complete
iteration (forward, forcing, adjoint, monthly update) before the year-long run.

% =============================================================================
\section{Outlook: a correlated background and gradient regularisation}
\label{sec:outlook}
% =============================================================================

The two principal challenges---gradient noise and an ineffective background---have
a common remedy. Replacing the diagonal background-error covariance $\bm{B}$ with
a spatially correlated one, and reparameterising the control vector through the
transform $\sigma_m = 1 + \bm{B}^{1/2}\bm{w}_m$ (a control-variable transform with
a Gaussian smoother as $\bm{B}^{1/2}$), accomplishes two things at once. It gives
the prior a meaningful spatial correlation length, so that information from
observations is spread physically rather than deposited in isolated cells; and,
because the gradient in the transformed variable is $\bm{B}^{1/2}$ applied to the
physical gradient, it low-pass filters the noisy adjoint gradient as an intrinsic
part of the preconditioning. A correlated observation-error model, or an
observation-error inflation to account for correlated retrieval errors, is a
complementary step. Calibrating the background correlation length and amplitude
against the known OSSE truth is itself a deliverable of the experiment.

% =============================================================================
\section{Summary}
% =============================================================================

We have set out a monthly 4D-Var system for terrestrial CO$_2$ flux estimation
built on the GCHP adjoint, evaluated in an OSSE. The scientific novelty is the
monthly-resolved control vector, which requires the adjoint gradient to be formed
as a difference of the backward accumulator across month boundaries; this is
derived and verified in the companion note. A single-month proof of concept
confirms that the assimilation loop functions and reduces the cost, and it
exposes two obstacles---a chaotic, noise-dominated adjoint gradient (shown to be
a genuine property of the convective adjoint, not a bug) and an under-weighted
background---whose common remedy is a spatially correlated background implemented
through a control-variable transform. The monthly, full-year experiment, together
with the OSSE evaluation of month-by-month analysis error, is the next step.

\end{document}
